% ============================================================
% Question Bank: ch02-10 Compound Interest Introduction
% Topic Title: Compound Interest Introduction
% CCSS: Supplementary
% Questions: 27 (15 MC + 6 SA + 6 Graphical)
% ============================================================

% --- Multiple Choice Questions (15) ---

\begin{practiceQuestion}{2-10-q01}{mc}
\begin{questionText}
Which statement best describes compound interest?
\end{questionText}
\choiceA{Interest is earned only on the original deposit}
\choiceB{Interest is earned on the principal plus any previously earned interest}
\choiceC{The interest rate changes every year}
\choiceD{Interest is only paid at the end of the investment}
\correctAnswer{B}
\explanation{Compound interest is ``interest on interest''---you earn interest on both the original principal and on any interest that has been added to the account.}
\end{practiceQuestion}

\begin{practiceQuestion}{2-10-q02}{mc}
\begin{questionText}
You deposit $\$600$ at $10\%$ annual compound interest. What is the balance after $1$ year?
\end{questionText}
\choiceA{$\$60$}
\choiceB{$\$610$}
\choiceC{$\$660$}
\choiceD{$\$700$}
\correctAnswer{C}
\explanation{After $1$ year: $600 \times 1.10 = \$660$.}
\end{practiceQuestion}

\begin{practiceQuestion}{2-10-q03}{mc}
\begin{questionText}
$\$800$ is deposited at $5\%$ annual compound interest. What is the balance after $2$ years?
\end{questionText}
\choiceA{$\$840$}
\choiceB{$\$880$}
\choiceC{$\$882$}
\choiceD{$\$900$}
\correctAnswer{C}
\explanation{Year 1: $800 \times 1.05 = \$840$. Year 2: $840 \times 1.05 = \$882$.}
\end{practiceQuestion}

\begin{practiceQuestion}{2-10-q04}{mc}
\begin{questionText}
In the formula $A = P(1 + r)^t$, what does $P$ represent?
\end{questionText}
\choiceA{The total amount after $t$ years}
\choiceB{The annual interest rate}
\choiceC{The original deposit (principal)}
\choiceD{The number of years}
\correctAnswer{C}
\explanation{$P$ stands for principal---the original amount deposited or invested.}
\end{practiceQuestion}

\begin{practiceQuestion}{2-10-q05}{mc}
\begin{questionText}
Compare: $\$1{,}000$ at $6\%$ simple interest for $2$ years vs.\ $\$1{,}000$ at $6\%$ compound interest for $2$ years. Which earns more and by how much?
\end{questionText}
\choiceA{Simple earns more by $\$3.60$}
\choiceB{Compound earns more by $\$3.60$}
\choiceC{They earn the same}
\choiceD{Compound earns more by $\$6.00$}
\correctAnswer{B}
\explanation{Simple: $1{,}000 \times 0.06 \times 2 = \$120$. Compound: Year 1: $1{,}060$. Year 2: $1{,}060 \times 1.06 = \$1{,}123.60$. Interest $= \$123.60$. Difference: $123.60 - 120 = \$3.60$.}
\end{practiceQuestion}

\begin{practiceQuestion}{2-10-q06}{mc}
\begin{questionText}
You invest $\$400$ at $10\%$ annual compound interest. What is the amount after $3$ years?
\end{questionText}
\choiceA{$\$520$}
\choiceB{$\$532.40$}
\choiceC{$\$540$}
\choiceD{$\$600$}
\correctAnswer{B}
\explanation{$A = 400(1.10)^3$. Year 1: $440$. Year 2: $484$. Year 3: $532.40$.}
\end{practiceQuestion}

\begin{practiceQuestion}{2-10-q07}{mc}
\begin{questionText}
Why does compound interest earn more than simple interest over the same period?
\end{questionText}
\choiceA{The rate is higher}
\choiceB{The principal changes}
\choiceC{Each year's interest is added to the balance, so the next year earns interest on a larger amount}
\choiceD{Simple interest uses a different formula}
\correctAnswer{C}
\explanation{With compound interest, the interest earned each period is added to the principal, increasing the base amount for future interest calculations.}
\end{practiceQuestion}

\begin{practiceQuestion}{2-10-q08}{mc}
\begin{questionText}
$\$2{,}000$ is invested at $4\%$ annual compound interest. What is the interest earned after $2$ years?
\end{questionText}
\choiceA{$\$160$}
\choiceB{$\$163.20$}
\choiceC{$\$200$}
\choiceD{$\$240$}
\correctAnswer{B}
\explanation{Year 1: $2{,}000 \times 1.04 = \$2{,}080$. Year 2: $2{,}080 \times 1.04 = \$2{,}163.20$. Interest: $2{,}163.20 - 2{,}000 = \$163.20$.}
\end{practiceQuestion}

\begin{practiceQuestion}{2-10-q09}{mc}
\begin{questionText}
A student says: ``$\$1{,}000$ at $5\%$ compound interest for $3$ years earns $\$150$---the same as simple interest.'' Is this correct?
\end{questionText}
\choiceA{Yes, both earn exactly $\$150$}
\choiceB{No, compound interest earns more than $\$150$}
\choiceC{No, compound interest earns less than $\$150$}
\choiceD{It depends on the bank}
\correctAnswer{B}
\explanation{Simple: $1{,}000 \times 0.05 \times 3 = \$150$. Compound: $1{,}000(1.05)^3 = \$1{,}157.625$, so interest $= \$157.63$. Compound earns about $\$7.63$ more.}
\end{practiceQuestion}

\begin{practiceQuestion}{2-10-q10}{mc}
\begin{questionText}
What is $(1 + 0.08)^2$?
\end{questionText}
\choiceA{$1.08$}
\choiceB{$1.16$}
\choiceC{$1.1664$}
\choiceD{$1.80$}
\correctAnswer{C}
\explanation{$(1.08)^2 = 1.08 \times 1.08 = 1.1664$.}
\end{practiceQuestion}

\begin{practiceQuestion}{2-10-q11}{mc}
\begin{questionText}
$\$300$ at $20\%$ compound interest for $2$ years gives a total of:
\end{questionText}
\choiceA{$\$360$}
\choiceB{$\$420$}
\choiceC{$\$432$}
\choiceD{$\$480$}
\correctAnswer{C}
\explanation{Year 1: $300 \times 1.20 = \$360$. Year 2: $360 \times 1.20 = \$432$.}
\end{practiceQuestion}

\begin{practiceQuestion}{2-10-q12}{mc}
\begin{questionText}
Which factor does NOT affect the amount of compound interest earned?
\end{questionText}
\choiceA{The principal}
\choiceB{The interest rate}
\choiceC{The number of years}
\choiceD{The color of the bank card}
\correctAnswer{D}
\explanation{Compound interest depends on the principal ($P$), the rate ($r$), and the time ($t$). The color of a card has no mathematical effect.}
\end{practiceQuestion}

\begin{practiceQuestion}{2-10-q13}{mc}
\begin{questionText}
$\$5{,}000$ is deposited at $3\%$ compound interest. What is the balance after $2$ years?
\end{questionText}
\choiceA{$\$5{,}150$}
\choiceB{$\$5{,}300$}
\choiceC{$\$5{,}304.50$}
\choiceD{$\$5{,}450$}
\correctAnswer{C}
\explanation{Year 1: $5{,}000 \times 1.03 = \$5{,}150$. Year 2: $5{,}150 \times 1.03 = \$5{,}304.50$.}
\end{practiceQuestion}

\begin{practiceQuestion}{2-10-q14}{mc}
\begin{questionText}
An account starts with $\$750$ at $8\%$ compound interest. After $1$ year the balance is $\$810$. How much interest was earned?
\end{questionText}
\choiceA{$\$8$}
\choiceB{$\$60$}
\choiceC{$\$80$}
\choiceD{$\$810$}
\correctAnswer{B}
\explanation{Interest: $810 - 750 = \$60$. Check: $750 \times 0.08 = \$60$. \checkmark}
\end{practiceQuestion}

\begin{practiceQuestion}{2-10-q15}{mc}
\begin{questionText}
$\$1{,}000$ at $10\%$: After $2$ years, simple interest gives $\$1{,}200$. What does compound interest give?
\end{questionText}
\choiceA{$\$1{,}200$}
\choiceB{$\$1{,}210$}
\choiceC{$\$1{,}220$}
\choiceD{$\$1{,}250$}
\correctAnswer{B}
\explanation{Compound: Year 1: $1{,}100$. Year 2: $1{,}100 \times 1.10 = \$1{,}210$. The extra $\$10$ is interest on interest.}
\end{practiceQuestion}

% --- Short Answer Questions (6) ---

\begin{practiceQuestion}{2-10-q16}{sa}
\begin{questionText}
Find the total amount after $2$ years for $\$1{,}500$ at $6\%$ annual compound interest.
\end{questionText}
\correctAnswer{$\$1{,}685.40$}
\explanation{Year 1: $1{,}500 \times 1.06 = \$1{,}590$. Year 2: $1{,}590 \times 1.06 = \$1{,}685.40$.}
\end{practiceQuestion}

\begin{practiceQuestion}{2-10-q17}{sa}
\begin{questionText}
How much interest does $\$200$ earn at $5\%$ compound interest after $3$ years?
\end{questionText}
\correctAnswer{$\$31.53$}
\explanation{$A = 200(1.05)^3 = 200 \times 1.157625 = \$231.525$. Interest: $231.53 - 200 = \$31.53$.}
\end{practiceQuestion}

\begin{practiceQuestion}{2-10-q18}{sa}
\begin{questionText}
$\$3{,}000$ at $4\%$ annual compound interest for $2$ years. What is the total?
\end{questionText}
\correctAnswer{$\$3{,}244.80$}
\explanation{Year 1: $3{,}000 \times 1.04 = \$3{,}120$. Year 2: $3{,}120 \times 1.04 = \$3{,}244.80$.}
\end{practiceQuestion}

\begin{practiceQuestion}{2-10-q19}{sa}
\begin{questionText}
Compare $\$500$ at $8\%$ for $2$ years: simple vs.\ compound. How much more does compound earn?
\end{questionText}
\correctAnswer{$\$3.20$}
\explanation{Simple: $500 \times 0.08 \times 2 = \$80$. Compound: Year 1: $540$. Year 2: $540 \times 1.08 = \$583.20$. Interest: $\$83.20$. Difference: $83.20 - 80 = \$3.20$.}
\end{practiceQuestion}

\begin{practiceQuestion}{2-10-q20}{sa}
\begin{questionText}
What is $(1.06)^2$?
\end{questionText}
\correctAnswer{$1.1236$}
\explanation{$1.06 \times 1.06 = 1.1236$.}
\end{practiceQuestion}

\begin{practiceQuestion}{2-10-q21}{sa}
\begin{questionText}
$\$10{,}000$ at $2\%$ compound interest for $2$ years. What is the total interest earned?
\end{questionText}
\correctAnswer{$\$404$}
\explanation{Year 1: $10{,}000 \times 1.02 = \$10{,}200$. Year 2: $10{,}200 \times 1.02 = \$10{,}404$. Interest: $10{,}404 - 10{,}000 = \$404$.}
\end{practiceQuestion}

% --- Graphical Questions (3 GMC + 3 GSA) ---

\begin{practiceQuestion}{2-10-q22}{gmc}
\begin{questionText}
The table below shows year-by-year balances for a $\$1{,}000$ deposit at $10\%$ compound interest. What is the balance at the end of Year $3$?

\begin{center}
\begin{tabular}{|c|c|c|}
\hline
\textbf{Year} & \textbf{Interest Earned} & \textbf{End Balance} \\
\hline
$0$ & --- & $\$1{,}000$ \\
\hline
$1$ & $\$100$ & $\$1{,}100$ \\
\hline
$2$ & $\$110$ & $\$1{,}210$ \\
\hline
$3$ & ? & ? \\
\hline
\end{tabular}
\end{center}
\end{questionText}
\choiceA{$\$1{,}300$}
\choiceB{$\$1{,}310$}
\choiceC{$\$1{,}331$}
\choiceD{$\$1{,}350$}
\correctAnswer{C}
\explanation{Year $3$ interest: $1{,}210 \times 0.10 = \$121$. Balance: $1{,}210 + 121 = \$1{,}331$.}
\end{practiceQuestion}

\begin{practiceQuestion}{2-10-q23}{gmc}
\begin{questionText}
The bar graph compares the total amount after $2$ years for simple and compound interest on a $\$2{,}000$ deposit at $5\%$. Which bar shows the compound interest total?

\begin{center}
\begin{tikzpicture}[scale=0.5]
  \draw[thick,->] (0,0) -- (0,6) node[above]{\small Total (\$)};
  \draw[thick,->] (0,0) -- (6,0);
  \fill[funBlue!60] (0.75,0) rectangle (2.25,4);
  \fill[funGreen!60] (3.25,0) rectangle (4.75,4.1);
  \node[below,font=\small] at (1.5,0) {A};
  \node[below,font=\small] at (4,0) {B};
  \foreach \y/\lab in {2/2100,4/2200,4.1/2205} {
    \draw (-0.1,\y) -- (0.1,\y);
  }
  \node[left,font=\small] at (-0.1,4) {$\$2{,}200$};
  \node[above,font=\small] at (1.5,4) {$\$2{,}200$};
  \node[above,font=\small] at (4,4.1) {$\$2{,}205$};
\end{tikzpicture}
\end{center}
\end{questionText}
\choiceA{Bar A ($\$2{,}200$)}
\choiceB{Bar B ($\$2{,}205$)}
\choiceC{Both are compound}
\choiceD{Neither is compound}
\correctAnswer{B}
\explanation{Simple: $2{,}000 + 200 = \$2{,}200$ (Bar A). Compound: $2{,}000(1.05)^2 = \$2{,}205$ (Bar B). Compound is higher.}
\end{practiceQuestion}

\begin{practiceQuestion}{2-10-q24}{gmc}
\begin{questionText}
The line graph shows the growth of $\$500$ at $10\%$ annual compound interest. What is the balance at Year $2$?

\begin{center}
\begin{tikzpicture}[scale=0.7]
  \draw[thick,->] (0,0) -- (5,0) node[right]{\small Year};
  \draw[thick,->] (0,0) -- (0,5) node[above]{\small Balance (\$)};
  \foreach \x/\lab in {0/0,1/1,2/2,3/3} {
    \draw (\x,-0.1) -- (\x,0.1) node[below=3pt,font=\small] {\lab};
  }
  \foreach \y/\lab in {1/200,2/400,2.5/500,3/600,3.5/700,4/800} {
    \draw (-0.1,\y) -- (0.1,\y) node[left,font=\small,xshift=-2pt] {\lab};
  }
  \draw[thick,funOrange,mark=*,mark size=2pt] plot coordinates {(0,2.5) (1,2.75) (2,3.025) (3,3.325)};
  \node[above right,font=\small,funOrange] at (0,2.5) {$\$500$};
  \node[above right,font=\small,funOrange] at (1,2.75) {$\$550$};
  \node[above right,font=\small,funOrange] at (2,3.025) {?};
  \node[above right,font=\small,funOrange] at (3,3.325) {$\$665.50$};
\end{tikzpicture}
\end{center}
\end{questionText}
\choiceA{$\$600$}
\choiceB{$\$605$}
\choiceC{$\$610$}
\choiceD{$\$620$}
\correctAnswer{B}
\explanation{Year 2: $550 \times 1.10 = \$605$.}
\end{practiceQuestion}

\begin{practiceQuestion}{2-10-q25}{gsa}
\begin{questionText}
Complete the compound interest table for $\$700$ at $10\%$ per year. Find the end-of-Year-$2$ balance.

\begin{center}
\begin{tabular}{|c|c|c|}
\hline
\textbf{Year} & \textbf{Interest} & \textbf{Balance} \\
\hline
$0$ & --- & $\$700$ \\
\hline
$1$ & $\$70$ & $\$770$ \\
\hline
$2$ & ? & ? \\
\hline
\end{tabular}
\end{center}
\end{questionText}
\correctAnswer{$\$847$}
\explanation{Year $2$ interest: $770 \times 0.10 = \$77$. Balance: $770 + 77 = \$847$.}
\end{practiceQuestion}

\begin{practiceQuestion}{2-10-q26}{gsa}
\begin{questionText}
The table shows simple and compound interest balances for $\$1{,}000$ at $5\%$. How much more does compound interest earn after Year $3$?

\begin{center}
\begin{tabular}{|c|c|c|}
\hline
\textbf{Year} & \textbf{Simple Balance} & \textbf{Compound Balance} \\
\hline
$0$ & $\$1{,}000$ & $\$1{,}000$ \\
\hline
$1$ & $\$1{,}050$ & $\$1{,}050$ \\
\hline
$2$ & $\$1{,}100$ & $\$1{,}102.50$ \\
\hline
$3$ & $\$1{,}150$ & ? \\
\hline
\end{tabular}
\end{center}
\end{questionText}
\correctAnswer{$\$7.63$}
\explanation{Compound Year 3: $1{,}102.50 \times 1.05 = \$1{,}157.625 \approx \$1{,}157.63$. Simple: $\$1{,}150$. Difference: $1{,}157.63 - 1{,}150 = \$7.63$.}
\end{practiceQuestion}

\begin{practiceQuestion}{2-10-q27}{gsa}
\begin{questionText}
The bar graph shows how much interest is earned each year on a $\$2{,}000$ deposit at $5\%$ compound interest. Find the interest earned in Year $3$.

\begin{center}
\begin{tikzpicture}[scale=0.6]
  \draw[thick,->] (0,0) -- (0,5) node[above]{\small Interest (\$)};
  \draw[thick,->] (0,0) -- (8,0);
  \fill[funBlue!60] (0.5,0) rectangle (2,2);
  \fill[funGreen!60] (2.5,0) rectangle (4,2.1);
  \fill[funOrange!60] (4.5,0) rectangle (6,0);
  \node[below,font=\small] at (1.25,0) {Yr 1};
  \node[below,font=\small] at (3.25,0) {Yr 2};
  \node[below,font=\small] at (5.25,0) {Yr 3};
  \foreach \y/\lab in {1/50,2/100,3/150,4/200} {
    \draw (-0.1,\y) -- (0.1,\y) node[left,font=\small,xshift=-2pt] {\lab};
  }
  \node[above,font=\small] at (1.25,2) {$\$100$};
  \node[above,font=\small] at (3.25,2.1) {$\$105$};
  \node[above,font=\small] at (5.25,0) {?};
\end{tikzpicture}
\end{center}
\end{questionText}
\correctAnswer{$\$110.25$}
\explanation{After Year 2, balance: $2{,}000 + 100 + 105 = \$2{,}205$. Year 3 interest: $2{,}205 \times 0.05 = \$110.25$.}
\end{practiceQuestion}
